\section{Listas Encadeadas}

Lista encadeada é uma implementação de conjunto dinâmico cujos membros, denominados
nós, nas suas formas mais simples, detêm um campo de chave e um ponteiro \texttt{proximo}
que aponta para o próximo membro no conjunto.
Numa implementação típica, a lista possui dois ponteiros especiais: \texttt{head} e
\texttt{tail}.
O \texttt{head} aponta para o primeiro elemento, enquanto o segundo para o último.
No entanto, para simplificar condições de contorno nas operações de inserção e remoção,
principalmente no início da lista, adiciona-se um nó chamado de \texttt{sentinela}, que
não possui chave válida, é apontado pelo \texttt{head} e cujo \texttt{proximo} aponta
para o primeiro membro efetivo da lista.

As listas encadeadas possuem diversas variações: listas simplesmente encadeadas,
duplamente encadeadas e circulares.
Também podemos contar com variações que incluem elementos de mais de um tipo.
De forma geral, assim como em vetores não ordenados, é possível realizar inserções e
remoções de um dado membro em tempo constante caso a ordenação não seja requisitada, e
buscas em tempo linear.

As listas encadeadas possuem a vantagem de serem mais flexíveis e permitirem inserções sem
necessariamente conhecer o total de elementos a serem comportados, diferentemente dos
vetores, que precisam ser realocados à medida que o número de membros cresce.
No entanto, os vetores fazem melhor aproveitamento da cache devido ao princípio de
localidade espacial.
No caso de existência de ordenação, a vantagem das listas encadeadas diminui mais,
pois uma busca num vetor ordenado pode ser executada em tempo logarítmico, enquanto
nas listas continua sendo realizado em tempo linear.

\subsection{Lista Simplesmente Encadeada}

Consiste na implementação mais simples para uma lista encadeada.
A lista mantém três atributos: \texttt{head}, \texttt{tail} e o número de elementos.
Uma busca nessa lista consiste numa pesquisa sequencial: percorremos cada elemento,
seguindo os ponteiros de \texttt{proximo} até obter a chave desejada.
É importante frisar que o que é retornado não é um ponteiro para o membro cuja chave
é aquela sendo pesquisada, mas sim um ponteiro para o membro anterior.
Isso porque um nó não conhece aquele que o antecede, apenas o que lhe sucede.

\begin{function}
\SetAlgoLined
\caption{Busca-Lista($L$, $k$)}
$w \gets L.\Head$\;
\Enqto{$w.\Prox \ne \nil$ \E $w.\Prox.\Chave \ne k$}{
    $w \gets w.\Prox$\;
}
\Se{$w.\Prox = \nil$}{
    \Retorna{$\nil$}\;
}
\Retorna{$w$}\;
\end{function}

A inserção é feita tomando-se como argumento o nó \texttt{prev} que antecederá o novo
membro do conjunto.
Neste momento, é feito o ajuste dos ponteiros de \texttt{proximo}, a fim de que o
\texttt{proximo} de \texttt{prev} aponte para o novo membro, e o \texttt{proximo} desse
último aponte para o antigo sucessor de \texttt{prev}.
Se \texttt{prev} era o \texttt{tail} da lista, então o novo nó será o novo \texttt{tail},
o que pode ser feito verificando-se se o \texttt{proximo} do novo é $\nil$.

\begin{function}[H]
\SetAlgoLined
\caption{Insere-Lista($L$, prev, $x$)}
$x.\Prox \gets $ prev.\Prox\;
prev.\Prox $\gets x$\;
\Se{$x.\Prox = \nil$}{
    $L.\Tail \gets x$\;
}
$L.\NumElem \gets L.\NumElem + 1$\;
\end{function}

A remoção toma como argumento o nó \texttt{prev} que precede aquele que será removido.
Nessa operação, o \texttt{proximo} de \texttt{prev} irá apontar para o \texttt{proximo}
do nó retirado.
Se o removido era o \texttt{tail}, então \texttt{prev} passa a ocupar essa função.

\begin{function}[H]
\SetAlgoLined
\caption{Remove-Lista($L$, prev)}
\Se{$\mathrm{prev}.\Prox = \nil$}{
    \Retorna\;
}
$x \gets \mathrm{prev}.\Prox$\;
prev.\Prox $\gets x.\Prox$\;
\Se{$\mathrm{prev}.\Prox = \nil$}{
    $L.\Tail \gets$ prev\;
}
$L.\NumElem \gets L.\NumElem - 1$\;
\end{function}

\subsection{Lista Duplamente Encadeada}

Essa variante funciona de maneira análoga à encadeada simples.
O adicional é que os membros agora possuem um ponteiro de \texttt{anterior}, que aponta
para o nó antecessor.
Nesse caso, a busca pode retornar um ponteiro para o próprio elemento que possui a chave
pesquisada.

\begin{function}
\SetAlgoLined
\caption{Busca-Lista-D($L$, $k$)}
$w \gets L.\Head.\Prox$\;
\Enqto{$w \ne \nil$ \E $w.\Chave \ne k$}{
    $w \gets w.\Prox$\;
}
\Se{$w = \nil$}{
    \Retorna{$\nil$}\;
}
\Retorna{$w$}\;
\end{function}

\begin{function}[H]
\SetAlgoLined
\caption{Insere-Lista-D($L$, prev, $x$)}
$x.\Prox \gets $ prev.\Prox\;
prev.\Prox $\gets x$\;
$x.\Ant \gets \mathrm{prev}$\;
\eSe{$x.\Prox = \nil$}{
    $L.\Tail \gets x$\;
}
{
    $x.\Prox.\Ant \gets x$\;
}
$L.\NumElem \gets L.\NumElem + 1$\;
\end{function}

\begin{function}[H]
\SetAlgoLined
\caption{Remove-Lista-D($L$, $x$)}
\Se{$x = L.\Head$}{
    \Retorna\;
}
$\mathrm{prev} \gets x.\Ant$\;
$\mathrm{prev}.\Prox \gets x.\Prox$\;
\eSe{$\mathrm{prev}.\Prox = \nil$}{
    $L.\Tail \gets$ prev\;
}
{
    $x.\Prox.\Ant \gets \mathrm{prev}$\;
}
$L.\NumElem \gets L.\NumElem - 1$\;
\end{function}

\subsection{Lista Circular Duplamente Encadeada}

A lista encadeada circular é aquela em que o último nó aponta para o primeiro em vez de
$\nil$, ou seja, o \texttt{head} da lista equivale a $L.\Tail.\Prox$.
Ao adicionar o ponteiro \texttt{anterior} a cada membro, obtemos, então, a lista
circular duplamente encadeada, que combina todos os elementos vistos até então sobre listas.
As operações de busca, inserção e remoção são análogas àquelas vistas para listas
duplamente encadeadas.

Iremos expor a implementação dessa variação usando o nó sentinela.
Com a lista vazia, os ponteiros \texttt{proximo} e \texttt{anterior} da sentinela apontam
para ela mesma.
À medida que as inserções ocorrem, esses ponteiros são ajustados para refletir a
estrutura circular.

A busca continua retornando um ponteiro para o nó que contém a chave desejada.
O que muda é a condição de parada quando a chave não existe.
Em vez de encerrar quando atingir $\nil$, o critério de parada passa a ser revisitar a
sentinela.

\begin{function}
\SetAlgoLined
\caption{Busca-Lista-DC($L$, $k$)}
$w \gets L.\Head.\Prox$\;
\Enqto{$w \ne L.\Head$ \E $w.\Chave \ne k$}{
    $w \gets w.\Prox$\;
}
\Se{$w = L.\Head$}{
    \Retorna{$\nil$}\;
}
\Retorna{$w$}\;
\end{function}

A estrutura circular elimina a condição de contorno de quando inserimos ou eliminamos no
final da lista.
Desse modo, as inserções e remoções funcionam da mesma forma independentemente de onde as
realizamos na lista.
A única operação que deve ser coibida é a remoção da própria sentinela.

\begin{function}[H]
\SetAlgoLined
\caption{Insere-Lista-DC($L$, prev, $x$)}
$x.\Prox \gets \mathrm{prev}.\Prox$ \;
prev.\Prox $\gets x$\;
$x.\Ant \gets \mathrm{prev}$\;
$x.\Prox.\Ant \gets x$\;
$L.\NumElem \gets L.\NumElem + 1$\;
\end{function}

\begin{function}[H]
\SetAlgoLined
\caption{Remove-Lista-DC($L$, $x$)}
\Se{$x = L.\Head$}{
    \Retorna\;
}
$\mathrm{prev} \gets x.\Ant$\;
$\mathrm{prev}.\Prox \gets x.\Prox$\;
$x.\Prox.\Ant \gets \mathrm{prev}$\;
$L.\NumElem \gets L.\NumElem - 1$\;
\end{function}
